\documentclass[12pt,a4paper]{ctexart}

\usepackage{amsmath,amssymb}
\usepackage{geometry}
\usepackage{enumitem}
\usepackage{hyperref}
\usepackage{booktabs}
\usepackage{listings}

\geometry{left=2.5cm,right=2.5cm,top=2.5cm,bottom=2.5cm}
\setlist{nosep,leftmargin=2em}

\title{第5章\quad 质点的动量矩与动量矩守恒}
\author{}
\date{}

\begin{document}

\maketitle

\section{知识点总结}

\subsection{力矩}

\begin{itemize}
    \item \textbf{单个质点}：用 $\vec{F}$、$E_k$、$\vec{P}$ 描述动力学性质；但要讨论质点绕某定点的转动，需引入力矩。
    \item \textbf{力矩定义}（对定点 $O$）：$\vec{M}_O = \vec{r}\times\vec{F}$，大小 $|\vec{M}_O| = rF\sin\alpha$（$\alpha$ 为 $\vec{r}$ 与 $\vec{F}$ 夹角）。
    \begin{itemize}
        \item 方向：由右螺旋法则确定，垂直于 $\vec{r}$ 与 $\vec{F}$ 所成平面；
        \item 空间矢量：力矩是空间矢量，不局限于固定平面；
        \item 等价表述：$|\vec{M}_O| = \;($定点到力作用线的距离$)\;\times F = \;($定点到质点的距离$)\;\times F_\tau$。
    \end{itemize}
    \item \textbf{质点圆周运动下的力矩}（以圆心 $O$ 为定点）：
    \begin{itemize}
        \item 切向加速度 $a_\tau = \beta r$（$\beta$ 为角加速度），法向加速度 $a_n = \omega^2 r$；
        \item 力矩大小 $M_O = r m a_\tau = r^2 m\beta$。
    \end{itemize}
    \item \textbf{物理意义}：力矩是引起质点绕某一定点转动状态变化（获得角加速度）的原因。
\end{itemize}

\subsection{质点的动量矩}

\begin{itemize}
    \item \textbf{动量矩（角动量）定义}（对 $O$ 点）：
    \[
        \vec{L}_O = \vec{r}\times\vec{P} = \vec{r}\times m\vec{v},
    \]
    \begin{itemize}
        \item 大小 $L_O = rp\sin\varphi = mrv\sin\varphi$；
        \item 方向由右螺旋法则确定；
        \item \textbf{特例：圆周运动} $L = rp = mrv$（$\vec{v}\perp\vec{r}$）；
        \item 单位：$\mathrm{kg\cdot m^2/s}$。
    \end{itemize}
    \item 质点对某点的动量矩在过该点的任意轴上的投影就等于质点对该轴的动量矩。
    \item \textbf{平面运动}：对运动平面内某参考点 $O$ 的动量矩也称为对过 $O$ 垂直于运动平面的轴的动量矩。
    \item \textbf{例}：圆锥摆（$A$ 为摆球圆周轨道圆心，$B$ 为悬挂点，杆与竖直方向夹角 $\alpha$）。
    \begin{itemize}
        \item 对 $A$ 点：$\vec{L}_A = \vec{r}\times m\vec{v}$，$|\vec{L}_A| = mr^2\omega$，方向沿 $z$ 轴正方向（大小方向均不变）；
        \item 对 $B$ 点：$|\vec{L}_B| = mr^2\omega/\sin\alpha$，\textbf{方向不断变化}；$\vec{L}_B$ 沿 $z$ 轴投影 $L_{Bz} = L_B\sin\alpha = mr^2\omega$（等于对 $A$ 点的动量矩）。
    \end{itemize}
    \item \textbf{说明}：动量矩只与"垂距"有关——参考点到速度方向直线的垂直距离，不是参考点到质点的距离。
    \item \textbf{质点动量矩定理}（微分形式）：
    \[
        \vec{M} = \frac{\mathrm{d}\vec{L}}{\mathrm{d}t},\quad \vec{M}\mathrm{d}t = \mathrm{d}\vec{L}.
    \]
    推导：$\dfrac{\mathrm{d}\vec{L}}{\mathrm{d}t} = \dfrac{\mathrm{d}}{\mathrm{d}t}(\vec{r}\times m\vec{v}) = \vec{r}\times\dfrac{\mathrm{d}(m\vec{v})}{\mathrm{d}t} + \dfrac{\mathrm{d}\vec{r}}{\mathrm{d}t}\times m\vec{v} = \vec{r}\times\vec{F} + \vec{v}\times m\vec{v} = \vec{M}$。
    \item \textbf{质点动量矩定理}（积分形式）：
    \[
        \int_{t_1}^{t_2}\vec{M}\cdot\mathrm{d}t = \vec{L}_2 - \vec{L}_1.
    \]
    \begin{itemize}
        \item 物理意义：质点所受合力矩的冲量矩等于质点的动量矩的增量；
        \item 冲量矩是动量矩变化的原因，动量矩的变化是力矩对时间的累积结果。
    \end{itemize}
    \item \textbf{质点动量矩守恒定律}：当 $\sum \vec{M}_i = 0$（合外力矩为零）时，$\vec{L} = \text{常矢量}$。
    \begin{itemize}
        \item \textbf{物理本质}：与空间各向同性（空间旋转对称性）相联系；
        \item \textbf{适用条件}：系统合外力矩为零；适用于宏观、微观、高速、低速全范围；
        \item \textbf{特例}：质点只受有心力作用时，质点被限制在与 $\vec{L}$ 垂直的平面内运动，对力心动量矩守恒；
        \item \textbf{应用}：开普勒第二定律（行星在相等时间内扫过相等面积）由万有引力下的动量矩守恒推出。
    \end{itemize}
\end{itemize}

\subsection{质心系的动量矩定律}

\begin{itemize}
    \item \textbf{重要结论}：在质心系中不论质心系是惯性系还是非惯性系，动量矩定理仍适用。
    \begin{itemize}
        \item 原因：质心系中惯性力 $-m_i\vec{a}_C$ 对质心 $C$ 的力矩
        \[
            \vec{M}_{C\text{惯}} = \sum \vec{r}_{Ci}\times(-m_i\vec{a}_C) = -\biggl(\sum m_i\vec{r}_{Ci}\biggr)\times\vec{a}_C = 0.
        \]
        \item 因为 $\sum m_i\vec{r}_{Ci} = 0$（$\vec{r}_{Ci}$ 是相对质心的位矢）。
    \end{itemize}
    \item \textbf{质心系动量矩定理}：
    \[
        \vec{M}_C = \frac{\mathrm{d}\vec{L}_C}{\mathrm{d}t}.
    \]
    \item \textbf{体系动量矩分解}：$\vec{L} = \vec{L}_C + \vec{L}_{CM}$
    \begin{itemize}
        \item 轨道角动量 $\vec{L}_C = \vec{r}_C \times m_C\vec{v}_C$（质心绕定点的动量矩）；
        \item 固有角动量 $\vec{L}_{CM} = \sum \vec{r}_{Ci}\times m_i\vec{v}_{Ci}$（各质点相对质心运动的动量矩）。
    \end{itemize}
    \item \textbf{牛顿质点力学的体系}（中心公式 $\vec{F} = \dfrac{\mathrm{d}(m\vec{v})}{\mathrm{d}t}$ 的三种推广）：
    \begin{itemize}
        \item 乘 $\mathrm{d}t$ 积分：动量定理；
        \item 标积 $\mathrm{d}\vec{s}$ 积分：动能定理；
        \item 矢乘 $\vec{r}$：动量矩定理 $\vec{r}\times\vec{F} = \dfrac{\mathrm{d}(\vec{r}\times m\vec{v})}{\mathrm{d}t}$，$\vec{M} = \dfrac{\mathrm{d}\vec{L}}{\mathrm{d}t}$。
    \end{itemize}
\end{itemize}

\section{公式汇总}

\begin{enumerate}[leftmargin=3em,label=\textbf{公式\arabic*}：,ref=公式\arabic*]

    \item $\vec{M}_O = \vec{r}\times\vec{F}$，$|\vec{M}_O| = rF\sin\alpha$

    说明：力对定点 $O$ 的力矩；空间矢量，方向由右螺旋法则确定；$\alpha$ 为 $\vec{r}$ 与 $\vec{F}$ 夹角；既等于力臂乘以 $F$，也等于 $r$ 乘以力的切向分量。

    \item $M_O = rma_\tau = r^2 m\beta$（圆周运动）

    说明：圆周运动下对圆心的力矩大小；$\beta = \mathrm{d}\omega/\mathrm{d}t$ 为角加速度；与转动惯量 $mr^2$ 配合得到 $M = I\beta$。

    \item $\vec{L}_O = \vec{r}\times\vec{P} = \vec{r}\times m\vec{v}$

    说明：质点对 $O$ 点的动量矩（角动量）；方向由右螺旋法则确定；$|\vec{L}_O| = mrv\sin\varphi$；单位 $\mathrm{kg\cdot m^2/s}$；圆周运动 $L = mrv$。

    \item $\vec{M} = \dfrac{\mathrm{d}\vec{L}}{\mathrm{d}t}$，$\vec{M}\mathrm{d}t = \mathrm{d}\vec{L}$

    说明：质点动量矩定理（微分形式）；合力矩等于质点动量矩对时间的变化率；推导基于 $\vec{r}\times\dfrac{\mathrm{d}(m\vec{v})}{\mathrm{d}t}$。

    \item $\displaystyle\int_{t_1}^{t_2}\vec{M}\cdot\mathrm{d}t = \vec{L}_2 - \vec{L}_1$

    说明：质点动量矩定理（积分形式）；合力矩的冲量矩等于质点动量矩的增量。

    \item $\sum \vec{M}_i = 0 \Rightarrow \vec{L} = \text{常矢量}$

    说明：质点动量矩守恒定律；合外力矩为零时总动量矩守恒；与空间旋转对称性相联系；是自然界最普适的定律之一。

    \item $\dfrac{\mathrm{d}\vec{S}}{\mathrm{d}t} = \dfrac{L}{2m}$（面积速度）

    说明：开普勒第二定律的数学表述；面积速度 $\mathrm{d}\vec{S}/\mathrm{d}t$ 是常矢量（有心力场中对力心）。

    \item $\vec{M}_{C\text{惯}} = -\biggl(\sum m_i\vec{r}_{Ci}\biggr)\times\vec{a}_C = 0$

    说明：质心系中惯性力对质心的合力矩为零；因为 $\sum m_i\vec{r}_{Ci} = 0$（相对质心位矢的带权和为零）；这是动量矩定理在质心系中仍成立的根据。

    \item $\vec{M}_C = \dfrac{\mathrm{d}\vec{L}_C}{\mathrm{d}t}$（质心系动量矩定理）

    说明：不论质心系是惯性系还是非惯性系，质心系中动量矩定理仍成立；与牛顿定律在非惯性系中加惯性力即可使用类比。

    \item $\vec{L} = \vec{L}_C + \vec{L}_{CM}$，其中 $\vec{L}_C = \vec{r}_C\times m_C\vec{v}_C$，$\vec{L}_{CM} = \sum \vec{r}_{Ci}\times m_i\vec{v}_{Ci}$

    说明：体系动量矩分解为轨道角动量（质心绕定点运动）和固有角动量（各质点相对质心运动）之和；如地球的动量矩 = 自转动量矩 + 公转轨道动量矩。

\end{enumerate}

\end{document}
